\backupslide{A different object with a cyclic unit group}{
For odd primes $p$ and $n\ge1$,
\[
\key{(\Z/p^n\Z)^\times\text{ is cyclic of order }p^{n-1}(p-1).}
\]
For $n>1$, the ring is not a field: $p\cdot p^{n-1}=0$.
\[
\begin{array}{c|c|c}
&\F_{p^n}&\Z/p^n\Z\\\hline
\text{Units}&p^n-1&p^{n-1}(p-1)\\
\text{Characteristic}&p&p^n
\end{array}
\]
}{
This is additional material, not an application of the finite-field
theorem to a residue ring. Units are exactly the classes not divisible
by $p$, giving $p^n-p^{n-1}$. Both factors in the displayed zero product
are nonzero when $n>1$. The next slides give a primitive-root lifting
proof. The field comparison assumes a field of order $p^n$; we have
not proved existence for every prime power in the main lectures.
}

\backupslide{Two arithmetics with nine elements}{
In $\Z/9\Z$, the powers of $2$ are
\[
1,2,4,8,7,5,\qquad2^6=1.
\]
These are all six units; the nonzero residues $3,6$ are not units.
\spacer
In $\F_9$, all eight nonzero elements are units,
and a primitive element has order $8$.
}{
Calculate successive powers modulo $9$. The nonzero residue $3$
satisfies $3^2=0$. A field of nine elements can instead be constructed
by adjoining a square root of $2$ to $\F_3$, using Lecture 2's method.
This contrasts a six-element unit group with an eight-element
multiplicative field group, despite equal sizes of the underlying sets.
}

\backupslide{First lift a primitive root to modulo $p^2$}{
Let $p$ be odd, and let $g$ generate $\F_p^\times$.

Choose $h=g$ or $h=g+p$ so that
\[
\key{h^{p-1}\not\equiv1\pmod{p^2}.}
\]
This is possible because
\[
(g+p)^{p-1}\equiv g^{p-1}+p(p-1)g^{p-2}\pmod{p^2},
\]
and the added term is nonzero modulo $p^2$.
}{
Choose any integer representative $g$ prime to $p$. The two choices
have the same primitive reduction modulo $p$. The linear binomial
term is divisible by $p$ but not $p^2$; all remaining terms are
divisible by $p^2$. Thus both powers cannot equal 1 modulo $p^2$.
Write $h^{p-1}=1+pc$ with $p\nmid c$. It would be false to say
every primitive representative automatically lifts without adjustment.
}

\backupslide{Each further lift contributes a factor of $p$}{
For odd $p$, if $u=1+p^rc$ with $r\ge1$ and $p\nmid c$, then
\[
u^p=1+p^{r+1}c',\qquad p\nmid c'.
\]
The first binomial term is $p^{r+1}c$; all later terms
are divisible by $p^{r+2}$.
\spacer
Starting with $h^{p-1}=1+pc$ gives
\[
\key{v_p\bigl(h^{(p-1)p^k}-1\bigr)=k+1.}
\]
}{
Define $v_p(m)$ as the exponent of $p$ dividing the nonzero integer
$m$, if needed. For $2\le j<p$, the binomial coefficient contributes
$p$, so the $j$th term has valuation at least $1+rj\ge r+2$.
The last term has valuation $rp\ge r+2$ since $p\ge3$.
After division by $p^{r+1}$, the result is $c$ modulo $p$, so
it remains nonzero. Induction gives the displayed formula.
This is a direct special-case proof, not an invocation of a general
lifting-the-exponent lemma.
}

\backupslide{The lifted element has full order}{
Modulo $p^n$, let $d=\ord(h)$.
\[
p-1\mid d,\qquad d\mid(p-1)p^{n-1}.
\]
Thus $d=(p-1)p^j$ for some $0\le j\le n-1$.

The lifting calculation shows the least such $j$ is $n-1$.
\[
\key{d=(p-1)p^{n-1}.}
\]
So $h$ generates every unit modulo $p^n$.
}{
Reduction modulo $p$ forces $p-1$ to divide the order; Lagrange
in the unit group gives the upper divisibility. The valuation formula
says $h^{(p-1)p^j}\equiv1\pmod{p^n}$ exactly when $j+1\ge n$.
This proves the exact order, with a single adjusted integer $h$
working for every $n$. The case $n=1$ is the finite-field theorem.
}

\backupslide{The exception starts at eight}{
\[
(\Z/2\Z)^\times=\{1\},\qquad(\Z/4\Z)^\times=\{1,3\}.
\]
Both are cyclic. Modulo $8$, every unit squares to $1$:
\[
1^2=3^2=5^2=7^2=1\pmod8.
\]
There is no element of order $4$.

For $n\ge3$,
\[
\key{(\Z/2^n\Z)^\times\cong C_2\times C_{2^{n-2}}.}
\]
}{
$C_m$ denotes a cyclic group of order $m$. The direct-product
generators are $-1$ and $5$. Repeated squaring gives
$v_2(5^{2^k}-1)=k+2$, starting with $5-1=4$:
if $u=1+2^rc$, $r\ge2$, $c$ odd, then
$u^2-1=2^{r+1}c(1+2^{r-1}c)$ and the last factor is odd.
So 5 has order $2^{n-2}$, and its powers are exactly the residues
equal to 1 modulo 4. Every odd residue is in this subgroup or its
negative, and $-1$ is outside it. This proves the direct product.
Its largest element order is $2^{n-2}$, smaller than its size
$2^{n-1}$, so it is not cyclic. The modulo-8 example alone suffices
for a short classroom comparison.
}

\backupslide{Why the field proof does not apply to residue rings}{
In $\Z/8\Z$, the degree-two polynomial
\[
X^2-1
\]
has four roots: $1,3,5,7$.
\spacer
\key{The polynomial root bound required a field.}

An abelian unit group need not be cyclic.
}{
The finite-field proof chose an element failing a low-degree polynomial
equation. With zero divisors, cancellation fails in the factor-theorem
argument. Thus the residue-ring theorem needs an additional lifting
argument, and the theorem fails for powers of 2 starting at 8.
This returns to Lecture 1's discussion of zero divisors.
}
